% =====================================================================
\section{Адресные светодиоды WS2812}
% =====================================================================
\lessonintro
  {управлять лентой адресных светодиодов WS2812 по одному проводу и безопасно рассчитывать её питание}
  {урок 1 (неблокирующая работа по \code{millis()}), урок 4 (гамма-коррекция), урок 16 (отдельное питание)}
  {на ленте из 8 пикселей идут радуга и «бегущий огонь», а плата остаётся свободной}
  {лента, кольцо или модуль WS2812B на 8 пикселей; резистор \SI{330}{\ohm}; конденсатор \SI{1000}{\micro F}; для длинных лент --- блок питания \SI{5}{V}}

\subsection{Сначала: цепочка светодиодов, которые слушают по очереди}

Обычная RGB-лента требует по проводу на каждый цвет, и все её светодиоды горят одинаково. В ленте WS2812
всё иначе: в каждом пикселе есть крошечная микросхема. Вы шлёте по одному проводу длинную очередь цветов.
Первый пиксель забирает первые три числа (красный, зелёный, синий), остальное передаёт дальше, второй
забирает свои три --- и так по цепочке. В итоге каждый пиксель горит своим цветом, а плате нужен всего
один вывод.

\subsection{Протокол: нули и единицы разной длины}

Данные передаются одним проводом без отдельного такта. Каждый бит --- это импульс HIGH, за которым следует
LOW. Бит определяет \emph{длительность HIGH}: короткая --- «0», длинная --- «1». Весь бит занимает
около \SI{1.25}{\micro s}, то есть скорость примерно \SI{800}{kHz}.

\begin{figure}[H]
\centering
\begin{tikzpicture}[x=1.7cm,y=0.9cm,font=\sffamily\scriptsize]
  \node[left,text=accent] at (-0.2,3.5) {бит «0»};
  \draw[very thick,accent] (0,3) -- (0.5,3) -- (0.5,4) -- (1.5,4) -- (1.5,3) -- (3,3);
  \node[above] at (1,4.05) {HIGH $\approx\SI{0.4}{\micro s}$};
  \node[below,inner sep=2pt] at (2.1,2.9) {LOW $\approx\SI{0.85}{\micro s}$};
  \node[left,text=esp] at (-0.2,0.8) {бит «1»};
  \draw[very thick,esp] (0,0) -- (0.5,0) -- (0.5,1) -- (2.0,1) -- (2.0,0) -- (3,0);
  \node[above] at (1.25,1.05) {HIGH $\approx\SI{0.8}{\micro s}$};
  \node[below,inner sep=2pt] at (2.35,-0.1) {LOW $\approx\SI{0.45}{\micro s}$};
  \draw[dashed,gray] (0.5,-0.9) -- (0.5,4.4);
  \draw[dashed,gray] (3,-0.9) -- (3,4.4);
  \draw[<->,>=Stealth] (0.5,-0.7) -- node[below]{один бит $\approx\SI{1.25}{\micro s}$} (3,-0.7);
\end{tikzpicture}
\caption{Кодирование битов в WS2812. Значения ориентировочные: допуск около $\pm\SI{0.15}{\micro s}$,
точные цифры смотрите в datasheet своей микросхемы}
\end{figure}

На каждый пиксель приходится 24 бита: по 8 бит на зелёный, красный и синий (именно в этом порядке у
большинства лент). После всех данных нужна пауза LOW не короче $\approx\SI{50}{\micro s}$ --- сигнал сброса. Он
говорит ленте «кадр закончился, показывайте»: поэтому \code{strip.show()} вызывают после изменений.

\begin{figure}[H]
\centering
\begin{tikzpicture}[font=\sffamily\scriptsize,
  px/.style={blk,minimum width=2.6cm,minimum height=1cm,font=\sffamily\scriptsize}]
  \node[px,fill=blkper] (a) at (0,0) {пиксель 1};
  \node[px,fill=blkper] (b) at (3.9,0) {пиксель 2};
  \node[px,fill=blkper] (c) at (7.8,0) {пиксель 3};
  \node[blk,fill=blkrtc,minimum width=2.6cm,minimum height=1cm] (r) at (11.7,0) {пауза, сброс\\ $>\SI{50}{\micro s}$};
  \node[below=5pt of a,align=center] {G7\ldots G0\;R7\ldots R0\;B7\ldots B0\\ 24 бита};
  \draw[arr] (-2.3,0) node[left]{DIN} -- (a.west);
  \draw[arr] (a.east) -- node[above=2pt]{DOUT} (b.west);
  \draw[arr] (b.east) -- (c.west);
  \draw[arr] (c.east) -- (r.west);
\end{tikzpicture}
\caption{Кадр данных. Каждый пиксель отрезает первые 24 бита, а остальное передаёт дальше через DOUT}
\end{figure}

\subsection{Почему не \code{digitalWrite()}}

Импульсы длиной меньше микросекунды нужно формировать с точностью в десятки наносекунд. Вызов
\code{digitalWrite()} выполняется долго и неравномерно, а прерывания (Wi-Fi, USB, таймеры) могут вклиниться
посреди передачи, и цвета «поплывут». На ESP32-S3 для этого есть периферия \textbf{RMT} (Remote Control):
процессор записывает в её память список длительностей, а аппаратура выдаёт импульсы сама. Библиотека
\textbf{Adafruit NeoPixel} использует RMT, поэтому цвета остаются точными даже при работе Wi-Fi.

\subsection{Питание: считаем ток}

Каждый пиксель --- это три светодиода. Один пиксель на полной яркости белым цветом потребляет до
$\approx\SI{60}{mA}$ (3 цвета по \SI{20}{mA}). Лента из 30 пикселей --- уже до \SI{1.8}{A}. Верхняя оценка тока:
\[
  I \approx N \cdot \SI{60}{mA} \cdot \frac{\text{BRIGHTNESS}}{255},
\]
где $N$ --- число пикселей, а BRIGHTNESS --- значение в \code{setBrightness()}.

Для 8 пикселей на яркости 40 из 255 получаем около $8\cdot60\cdot0{,}16\approx\SI{75}{mA}$, и порт USB это
выдержит. Для 60 пикселей на полной яркости нужен блок питания \SI{5}{V} на \SI{3}{A}.

\begin{figure}[H]
\centering
\begin{circuitikz}[scale=0.95,transform shape]
  \node[blk,fill=blkcpu,minimum width=2.8cm,minimum height=4cm] (m) at (0,0) {};
  \node[blk,fill=blkper,minimum width=2.8cm,minimum height=4cm] (s) at (9,0) {};
  \node[font=\sffamily\small\bfseries] at (0,2.4) {ESP32-S3};
  \node[font=\sffamily\small\bfseries,align=center] at (9,2.4) {Лента WS2812};
  \node[font=\sffamily\scriptsize,anchor=east] at (1.4,1.5) {\pin{16}};
  \node[font=\sffamily\scriptsize,anchor=east] at (1.4,-1.5) {GND};
  \node[font=\sffamily\scriptsize,anchor=west] at (7.7,1.5) {DIN};
  \node[font=\sffamily\scriptsize,anchor=west] at (7.7,0) {+5 В};
  \node[font=\sffamily\scriptsize,anchor=west] at (7.7,-1.5) {GND};
  \draw[thick,okgreen] (1.4,1.5) to[R,l=\SI{330}{\ohm}] (7.6,1.5);
  \draw[thick,black!80] (1.4,-1.5) -- (7.6,-1.5);
  \node[blk,fill=blkrtc,minimum width=2.4cm,minimum height=1cm,font=\sffamily\scriptsize] (ps) at (3.4,0) {Блок питания\\ +5 В};
  \draw[thick,esp] (ps.east) -- (7.6,0);
  \draw[thick,black!80] (3.4,-0.5) -- (3.4,-1.5);
  \fill (3.4,-1.5) circle (1.6pt);
  \draw[thick,esp] (6.0,0) to[C,l_=\SI{1000}{\micro F}] (6.0,-1.5);
  \fill[esp] (6.0,0) circle (1.6pt);
  \fill (6.0,-1.5) circle (1.6pt);
  \node[font=\sffamily\scriptsize] at (4.5,-2.5) {общая земля: минус блока питания соединён с GND платы};
\end{circuitikz}
\caption{Подключение ленты. Резистор в линии данных гасит выбросы, конденсатор на питании
смягчает бросок тока при включении, а земля платы и блока питания общая}
\end{figure}

\begin{caution}
Сначала подключайте GND, затем питание, а сигнал --- в последнюю очередь. Подача данных на ленту без общей земли может
вывести из строя первый пиксель. Конденсатор должен быть рассчитан хотя бы на \SI{6.3}{V}, его
\textbf{полярность} (минус на полоске корпуса) нельзя путать.
\end{caution}

Резистор \SIrange{330}{470}{\ohm} ставьте как можно ближе к первому пикселю. Лента питается от
\SI{5}{V}, и по datasheet ей нужен вход выше $0{,}7\cdot\SI{5}{V}=\SI{3.5}{V}$, а плата даёт \SI{3.3}{V}. Часто это
работает, но на длинных проводах может сбоить. Надёжный вариант --- микросхема-преобразователь уровня
(например, 74AHCT125) в линии данных; она превращает \SI{3.3}{V} в \SI{5}{V}.

\subsection{Цвет через HSV}

Смешивать красный, зелёный и синий для радуги неудобно. Модель \textbf{HSV} описывает цвет понятнее:
\emph{Hue} (оттенок) --- положение на цветовом круге от 0° до 360°, \emph{Saturation} (насыщенность) ---
насколько цвет яркий, а не бледный, \emph{Value} (яркость). Чтобы получить радугу, достаточно менять один
Hue.

\begin{figure}[H]
\centering
\begin{tikzpicture}
  \foreach \h in {0,3,...,357}{
    \pgfmathsetmacro{\hh}{\h/360}
    \definecolor{hcol}{hsb}{\hh,1,1}
    \fill[hcol] ({\h/360*14},0) rectangle ({(\h+3)/360*14},0.8);
  }
  \draw (0,0) rectangle (14,0.8);
  \foreach \a/\n in {0/красный,60/жёлтый,120/зелёный,180/голубой,240/синий,300/пурпурный}
    \draw ({\a/360*14},0) -- ++(0,-0.15) node[below,font=\sffamily\scriptsize,align=center]{\n\\ \a°};
  \node[above,font=\sffamily\scriptsize] at (7,0.85) {Hue (оттенок): 0\ldots360°, потом круг замыкается};
\end{tikzpicture}
\caption{Шкала оттенков HSV. Ползунок Hue проходит все цвета радуги}
\end{figure}

\subsection{Встроенный светодиод как лента из одного пикселя}

На многих платах ESP32-S3-DevKitC-1 есть светодиод WS2812: на ревизии v1.0 он подключён к \pin{48}, а
на v1.1 --- к \pin{38} (мы использовали его в уроке 1). Это та же «лента», только из одного пикселя.
В ядре Arduino-ESP32 для неё есть готовая функция \code{rgbLedWrite(pin, r, g, b)}. Пример ниже без
всякой ленты показывает радугу из HSV.

\codecap{Радуга на встроенном светодиоде: HSV вручную}
\codeinput{l17_pixel_builtin}

\subsection{Эффекты без \code{delay()}}

В примере с лентой каждый эффект --- это функция одного шага: она сдвигает картинку и вызывает
\code{show()}. Шаги запускаются по таймеру \code{millis()}, поэтому \code{loop()} не блокируется и
плате остаётся время на кнопки, датчики и сеть (урок 5).
Радуга использует Hue пикселя $= \text{сдвиг} + i\cdot65536/N$; бегущий огонь гасит ленту и зажигает один
пиксель, сдвигая его на шаг.

Соберём программу в два шага.

\subsubsection{Шаг 1. Радуга}

Лента подключена по схеме выше. Радуга: оттенок каждого пикселя сдвинут относительно соседа, а весь
узор понемногу сдвигается со временем.

\codecap{Шаг 1. Радуга на ленте из 8 пикселей}
\codeinput{l17_strip_s1}

\begin{itemize}[leftmargin=*]
  \item \code{Adafruit\_NeoPixel strip(\ldots)} --- объект ленты: число пикселей, вывод и порядок цветов;
  \item \code{rainbowStep()} задаёт каждому пикселю свой Hue, а \code{hueOffset} сдвигает весь узор;
  \item шаг запускается по \code{millis()} каждые \SI{40}{ms}, а \code{show()} отправляет кадр на ленту.
\end{itemize}

\subsubsection{Шаг 2. Бегущий огонь и переключение эффектов}

Добавим второй эффект и будем переключать их по таймеру. Каждый эффект остаётся функцией одного шага,
поэтому \code{loop()} по-прежнему ничего не блокирует.

\codecap{Шаг 2. Радуга и бегущий огонь на ленте из 8 пикселей}
\codeinput{l17_strip}

\begin{itemize}[leftmargin=*]
  \item \code{chaseStep()} гасит ленту (\code{clear()}) и зажигает один пиксель \code{head}, сдвигая его на шаг;
  \item \code{rainbowMode} --- флаг текущего эффекта; он меняется каждые \SI{8}{s};
  \item в \code{loop()} выбирается нужная функция шага, а вызов \code{show()} остаётся один.
\end{itemize}

Установите библиотеку \textbf{Adafruit NeoPixel}. Функция \code{gamma32()} применяет ту же
гамма-коррекцию, что и урок 4: без неё середина радуги выглядит блёклой.

\subsection{Самостоятельный опыт: шкала уровня}

Сделайте отдельный скетч, в котором потенциометр на \pin{1} задаёт, сколько пикселей горит: от 0 до 8. Первые пиксели
зелёные, средние жёлтые, последний красный. Подсказки: \code{map(analogRead(...), 0, 4095, 0, NUM\_PIXELS + 1)}
и перебор пикселей циклом \code{for}. Обновляйте ленту только когда число пикселей изменилось, а не на
каждом проходе \code{loop()}.

\subsection{То же на ESPUI: выбираем цвет ползунком}

Радугу из HSV мы раньше только смотрели, а теперь можно выбирать цвет самому. Два ползунка в
панели --- «Hue» от 0° до 359° и «Яркость» --- превращаются в тот же перевод HSV~$\to$~RGB,
который уже есть в этом уроке. Никакой ленты не нужно: используем встроенный светодиод.
Программу соберём в три шага.

\subsubsection{Шаг 1. HSV в RGB без сети}

Сначала проверим сам перевод цвета. Функция \code{hsvToRgb()} превращает оттенок, насыщенность и
яркость в три числа для светодиода, а \code{loop()} просто крутит оттенок по кругу.

\codecap{Шаг 1. Радуга на встроенном светодиоде}
\codeinput[\footnotesize]{l17_espui_hue_s1}

\begin{itemize}[leftmargin=*]
  \item \code{hue / 60} --- номер одного из шести секторов цветового круга, остаток задаёт положение внутри сектора;
  \item \code{rgbLedWrite(pin, r, g, b)} зажигает встроенный светодиод нужным цветом;
  \item оттенок растёт на единицу и замыкается по кругу: \code{(hue + 1) \% 360}.
\end{itemize}

\subsubsection{Шаг 2. Оттенок выбираем ползунком}

Вместо автоматического вращения подключим панель: ползунок «Hue» записывает оттенок в
\code{hueValue}, а \code{loop()} использует это значение.

\codecap{Шаг 2. Ползунок Hue в панели ESPUI}
\codeinput[\footnotesize]{l17_espui_hue_s2}

\begin{itemize}[leftmargin=*]
  \item \code{hueCallback} только записывает число в \code{std::atomic}, светодиодом управляет \code{loop()};
  \item \code{ESPUI.slider(\ldots, 0, 0, 359)} --- ползунок оттенка: начальное значение 0, диапазон 0--359;
  \item яркость пока зашита в вызов \code{hsvToRgb(\ldots, 40, \ldots)}.
\end{itemize}

\subsubsection{Шаг 3. Яркость и подпись с RGB}

Добавим второй ползунок и подпись с итоговыми числами RGB.

\codecap{Шаг 3. ESPUI: оттенок и яркость встроенного RGB}
\codeinput[\footnotesize]{l17_espui_hue}

\begin{itemize}[leftmargin=*]
  \item второй ползунок пишет яркость 0--255 в \code{brightness}; она передаётся в \code{hsvToRgb()} как \code{val};
  \item подпись \code{ESPUI.label(\ldots)} показывает цвет в виде «R, G, B»;
  \item подпись обновляется раз в \SI{300}{ms}, а светодиод --- на каждом проходе \code{loop()}.
\end{itemize}

Подключитесь к сети \code{S3-Hue} (пароль \code{lesson-39-demo}). Здесь режим
\code{sliderContinuous} включён намеренно: цвет должен меняться, пока вы двигаете палец. Если
плата и браузер не успевают, верните \code{false}: событие придёт после отпускания. На ревизии
v1.1 замените \code{RGB\_PIN} на 38.

\begin{trouble}
\begin{itemize}[leftmargin=*]
  \item \textbf{Лента не горит или горит случайными цветами} --- сигнал подан на DOUT вместо DIN, нет общей земли или не тот вывод в программе.
  \item \textbf{Цвета не те (вместо красного зелёный)} --- у ленты другой порядок цветов: замените \code{NEO\_GRB} на \code{NEO\_RGB} в конструкторе.
  \item \textbf{Плата перезагружается при яркости выше нескольких пикселей} --- просадка питания: снизьте \code{setBrightness()} или питайте ленту от блока \SI{5}{V}.
\end{itemize}
\end{trouble}

\begin{summary}
\begin{itemize}[leftmargin=*]
  \item WS2812 принимает по одному проводу 24 бита на пиксель; бит «0» и «1» отличаются длиной HIGH около \SI{0.4}{\micro s} и \SI{0.8}{\micro s}.
  \item Точные импульсы формирует RMT, поэтому используйте библиотеку, а не \code{digitalWrite()}; ток считайте как $N\cdot\SI{60}{mA}$ на полной яркости.
  \item Для эффектов используйте HSV и шаги по \code{millis()}, а питание ленты и общую землю соединяйте по правилам.
\end{itemize}
\end{summary}

\begin{selfcheck}
\begin{enumerate}
  \item Чем отличаются в протоколе WS2812 нулевой и единичный биты?
  \item Какой ток потребляет лента из 30 пикселей на полной яркости белого и что с этим делать?
  \item Зачем нужны резистор в линии данных и конденсатор на питании?
\end{enumerate}
\end{selfcheck}

\begin{tasks}
\begin{enumerate}
\item \lv{1} Зажгите первые три пикселя красным, зелёным и синим цветом.
\item \lv{1} Измените яркость всей ленты и посчитайте потребляемый ток для вашего числа светодиодов.
\item \lv{2} Сделайте «бегущий огонёк», который перемещается по ленте без \code{delay()}.
\item \lv{2} Измените скорость радуги потенциометром и направление движения кнопкой.
\item \lv{3} Покажите на ленте температуру DS18B20: синий --- холодно, зелёный --- комфортно, красный --- жарко.
\item \lv{4} Сделайте эффект «дыхание» с гамма-коррекцией, не используя \code{delay()}.
\end{enumerate}
\end{tasks}
